% !TEX TS-program = xelatex
% !TEX encoding = UTF-8 Unicode
% !Mode:: "TeX:UTF-8"

\documentclass{resume}
\usepackage{zh_CN-Adobefonts_external} % Simplified Chinese Support using external fonts (./fonts/zh_CN-Adobe/)
% \usepackage{NotoSansSC_external}
% \usepackage{NotoSerifCJKsc_external}
% \usepackage{zh_CN-Adobefonts_internal} % Simplified Chinese Support using system fonts
\usepackage{linespacing_fix} % disable extra space before next section
\usepackage{cite}

\begin{document}
\pagenumbering{gobble} % suppress displaying page number

\name{秦亚飞}

\basicInfo{
  \email{qinyafee@163.com} \textperiodcentered\
  \phone{(+86) 18501689646} \textperiodcentered\
  \faBirthdayCake{ 1990/07} \textperiodcentered\
  \faMapMarker{ 上海} \textperiodcentered}


\section{\faBook\ 自我评价}
% increase linespacing [parsep=0.5ex]
8年智驾开发/管理，量产经验丰富，自驱力强；熟悉BEV感知、数据合成、闭环仿真、自动标注、地图定位开发业务

\section{\faCogs\ 专业技能}
% increase linespacing [parsep=0.5ex]
\begin{itemize}[parsep=0.5ex]
  \item 语言：C++、Python、PyTorch
  \item 重建：熟悉3dgs、streetGaussian、driveStudio；driveDreamer；dreamfusion、dreamGaussian、Hunyuan
  \item 生成：熟悉VAE、Diffusion、DiT、DDPM、Flow matching基础理论，熟悉经典工作如ControlNet、LoRA、dreambooth、MagicDrive、openSora等
  \item 感知：针孔/鱼眼内参模型、图像/特征点矫正、BEV转换、各任务经典模型/数据集/评测方法
  \item slam：图优化、EKF/UKF、粒子滤波、预积分；JPDA关联、HM匹配、车道线匹配；开源框架：ORB\_SLAM/VINS、MSCKF、g2o、Ceres Solver
  \item 优达学院无人驾驶工程师纳米学位（2017/10--2018/8）
\end{itemize}

\section{\faGraduationCap\  教育背景}
\datedsubsection{\textbf{浙江大学} \qquad 能源工程学系 \qquad  能源与环境系统工程 \qquad 本科}{2010/8--2014/6}

\section{\faUsers\ 工作经历}
\datedsubsection{\textbf{路特斯科技有限公司} \quad 算法软件部 \quad  环境模型算法专家}{2022/7--今}
\begin{onehalfspacing}

  负责数据合成业务
  \begin{itemize}
    \item 主导感知训练数据合成、端到端demo、闭环仿真业务
    \item 带领团队开发预处理、3dgs场景重建、3d资产生成、在2d视频流/3d场景中插入资产、场景可控生成模型、闭环仿真框架
    \item 可重建行人等柔性物体；可生成360°资产，如海外TSI、锥桶、水马、异型车辆、异型障碍物等
    \item 目前已经上线，支持修改传感器内外参、动/静态障碍物编辑、红绿灯控制，支持NVS渲染及全链路闭环仿真
    \item 交付纯视觉点云，支持无lidar静态元素标注
    \item 场景可控生成：可生成多视角图片/短视频，支持天气、光照、障碍物等控制
  \end{itemize}

  负责泊车感知建图的开发
  \begin{itemize}
    \item 主导IPM拼接、停车位/BEV freespace/OD检测、后处理、记忆行/泊车建图、车位到车位等开发任务
    \item 制定感知数据闭环方案，包括研发数采计划、标注、回测等环节
    \item 支持4D标注，提供点云建图算子；动态障碍物已经实现自动标注，极大节省标注费用
  \end{itemize}
  负责全场景定位算法开发
  \begin{itemize}
    \item {搭建算法框架，带领团队开发定位模块：支持高速/城市NOA全局定位、行车/泊车局部定位；适配不同配置：Hdmap/SDmap、高精/低精imu、有/无电机转速}
    \item 制定定位SIL方案，真值生产、数据切片/清洗、KPI评测、VIZ可视优化，与上下游定义接口、KPI、状态预警
    \item 已在lotus全系国内/欧洲、领克Z10车型OTA量产
  \end{itemize}
  负责lambda高快NOA定位供应商管理
  \begin{itemize}
    \item {分析jira，与感知/规控/地图/传感器/系统等沟通，界定定位模块问题，并推动供应商制定优化方向，跟踪执行}
    \item 验收供应商定位KPI报告，review量产问题整改方案
  \end{itemize}
\end{onehalfspacing}

\datedsubsection{\textbf{小鹏汽车科技有限公司} \quad 地图开发部 \quad  高精地图算法工程师}{2021/5--2022/7}
\begin{onehalfspacing}
  高速高架更新图层开发
  \begin{itemize}
    \item {负责车端更新图层模块开发，实现云端数据下载和地图错误修复功能，已量产交付；熟悉ADASISV3协议}
    \item 负责云端质检流程开发，支持地图数据几何/逻辑问题检查
  \end{itemize}

  城市更新图层开发
  \begin{itemize}
    \item 负责车道线生成算法，利用量产车众包上报的感知定位等数据，构建全局地图，修复地图错误或缺失区域
  \end{itemize}
\end{onehalfspacing}

\datedsubsection{\textbf{哈曼（中国）投资有限公司} \quad ADAS BU \quad 计算机视觉算法工程师}{2018/9--2021/5}
\begin{onehalfspacing}
  参与APA开发，负责建图定位模块的开发
  \begin{itemize}
    \item 基于360环视系统，融合几何特征点、语义特征和轮速计信息，实现VWO视觉里程计
    \item 熟悉手眼标定原理，掌握里程计、环视系统标定方法
  \end{itemize}

  负责Camera/ LiDAR外参标定工具的开发
  \begin{itemize}
    \item 在线标定：实现基于边缘匹配的在线标定算法
    \item 离线标定：掌握基于AR tags的标定工具
  \end{itemize}

  参与TJA（交通拥堵辅助系统）的预研，负责融合跟踪模块的开发
  \begin{itemize}
    \item 基于Carla模拟生成激光雷达和毫米波雷达的目标检测列表，跟踪多目标轨迹
  \end{itemize}
\end{onehalfspacing}

\datedsubsection{\textbf{上海汽车集团商用车公司技术中心}  \quad  动力总成与电气化中心 \quad  系统工程师 }{2014/8--2018/9}
负责新能源关键系统的需求分析和技术开发；负责设计方案及变更发布、供应商管理；跨部门合作智驾技术预研

\begin{onehalfspacing}
  \begin{itemize}
    \item 开发电机及控制器、电子水泵/冷却风扇，配合整车ESP、防盗、远程监控、驻坡等功能的开发
    \item 推动解决各类NVH、EMC干扰、动力电池MSD击穿等复杂问题
    \item 园区接驳车智驾demo演示
  \end{itemize}
\end{onehalfspacing}

% Reference Test
%\datedsubsection{\textbf{Paper Title\cite{zaharia2012resilient}}}{May. 2015}
%An xxx optimized for xxx\cite{verma2015large}
%\begin{itemize}
%  \item main contribution
%\end{itemize}

\section{\faCode\ 其他}
\begin{itemize}[parsep=0.5ex]
  \item 多目VIO系统实现与动态初始化深度优化（VR头显设备）{2021/2--2022/5}
  \item CN117310756B 多传感器融合定位方法及系统、机器可读存储介质 2023.12.29
  \item CN117249839A 一种车道线地图生成方法、装置、设备和存储介质 2023.12.19
\end{itemize}

%% Reference
%\newpage
%\bibliographystyle{IEEETran}
%\bibliography{mycite}
\end{document}
